\chapter{Conclusion}

L'application couvre l'ensemble du cahier des charges : les 23 besoins fonctionnels (P1, P2 et P3), les
16 règles de gestion, contrôlées dans la couche service, et les besoins non fonctionnels de sécurité, de
fiabilité et d'ergonomie. La démarche du modèle en Y a permis de fixer les cas d'utilisation et la
conception avant le code, puis de faire évoluer les diagrammes de conception avec l'implémentation.
Le passage de commande, transactionnel et protégé contre les accès concurrents, constitue le cœur
technique du projet.
